%%%%%%%%% JPS abstract %%%%%%%%%%%%%%%%%%%%%%%%%%%%%%%%%%%%%%%%%%
\documentclass[12pt,a4paper]{jsarticle}

%%%%%%%%% packages %%%%%%%%%%%%%%%%%%%%%%%%%%%%%%%%%%%%%%%%%%%%%%
\usepackage{graphicx} % Include figure files
\usepackage[%                           % 余白の設定
%mag=1400,%                              jarticle の場合（14ptに）
dvipdfm,truedimen,%
top=30truemm,bottom=20truemm,%
left=20truemm,right=20truemm]{geometry}
\pagestyle{empty}

%%%%%%%%% header %%%%%%%%%%%%%%%%%%%%%%%%%%%%%%%%%%%%%%%%%%%%%%%%
\begin{document}
\vspace{-5pt}
\begin{center}
{\gt \Large Higgs challenge }\\[14pt]

{\gt \large NH m1 \\ 髙橋宏太 \\ 2026年2月27日}\\[5pt]



\end{center}

\vspace{10pt}
%%%%%%%%% main %%%%%%%%%%%%%%%%%%%%%%%%%%%%%%%%%%%%%%%%%%%%%%%%%%

\section{アプローチ}
基本的にGeminiのProの助けを借りた。
アーキテクチャの選択はkaggleのhiggs challengeで優勝した方の記述を参考にして、Geminiに生成してもらった。
アーキテクチャを決定したのちに、epoch数やlayer数を変えて予想のams scoreや過学習をチェックし、最大値を出す方向に修正していった。

\section{アーキテクチャ}
Geminiの提案で物理的な描像を踏まえた特徴量、観測された質量の比率（\texttt{NEW\_mass\_ratio}）、全粒子の横運動量の合計（\texttt{NEW\_total\_pt}）、ジェット数とエネルギーを掛け合わせたもの（\texttt{NEW\_jet\_energy\_sum}）を追加した。

構造はニューラルネットワーク構造を使った。
1層目は、Sparse Connectivityという各ニューロンが10個の入力しか見ないように制限した。
これにより、過学習を防ぎ、頑健な特徴を見つけさせた。
2層目はFully Connected(全結合層)として用いたが、Max channel活性化関数という、600個のニューロンを3個づつのグループに分け、3つの中で最も出力が大きかった一個だけ通過させる関数を使った。
これにより、非線形性の表現や、情報の競合による、識別能力の向上を図った。

交差検証、アンサンブル学習をおこなった。交差検証については、データを検証用と学習用の2個に分割した。
その後データの並び替えを3回行い、計6つの学習と検証のモデルを生成し、予想のams scoreを算出した。

\section{予想ams score}
予想ams scoreは3.5049, 閾値は0.8385だった。
\end{document}
